\chapter{Euclidean整环、主理想整环与唯一分解整环}

有许多类环比一般的环具有更多代数结构。本章考察三类这样的环：具有除法算法的环（Euclidean整环）、每个理想都是主理想的环（主理想整环），以及元素能分解为素元的环（唯一分解整环）。这类环的主要例子是整数环 $\ZZ$ 和系数在域 $F$ 中的多项式环 $F[x]$。在这里，我们把绪论中关于整数 $\ZZ$ 陈述过的全部定理证明为更一般环论结果的特殊情形。这些结果将在下一章应用到特殊情形 $F[x]$ 中。

本章中所有的环都是交换环。

\section{Euclidean整环}

我们首先定义整环 $R$ 上范数的概念。这里的范数本质上只是 $R$ 中元素“大小”的一种度量。

\noindent\textbf{定义.}
设 $R$ 是整环。任意满足 $N(0)=0$ 的函数
\[
N:R\longrightarrow \ZZ^+\cup\{0\}
\]
称为 $R$ 上的一个\emph{范数}。若对所有 $a\neq0$ 都有 $N(a)>0$，则称 $N$ 为\emph{正范数}。

注意，这里的“范数”概念相当弱；同一个整环 $R$ 可能有若干不同的范数。

\noindent\textbf{定义.}
整环 $R$ 称为\emph{Euclidean整环}（或称它具有\emph{除法算法}），若存在 $R$ 上的一个范数 $N$，使得对任意 $a,b\in R$，$b\neq0$，存在 $q,r\in R$ 满足
\[
a=qb+r,
\qquad r=0 \text{ 或 } N(r)<N(b).
\]
元素 $q$ 称为商，$r$ 称为余数。

除法算法的重要性在于：它允许我们对 $R$ 中两个元素 $a,b$ 进行 Euclidean 算法。通过连续“除法”（这些除法实际上发生在 $R$ 的分式域中），可以写成
\begin{align*}
	a&=q_0b+r_0,\tag{0}\\
	b&=q_1r_0+r_1,\tag{1}\\
	r_0&=q_2r_1+r_2,\tag{2}\\
	&\ \vdots\\
	r_{n-2}&=q_nr_{n-1}+r_n,\tag{$n$}\\
	r_{n-1}&=q_{n+1}r_n.\tag{$n+1$}
\end{align*}
其中 $r_n$ 是最后一个非零余数。这样的 $r_n$ 必然存在，因为只要余数非零，就有
\[
N(b)>N(r_0)>N(r_1)>\cdots>N(r_n),
\]
这是一列递减的非负整数，不可能无限继续。注意这些元素一般并不唯一。

\paragraph{例}
\begin{enumerate}[label=(\arabic*)]
	\setcounter{enumi}{-1}
	\item 域是 Euclidean整环的平凡例子：任意范数都满足条件，例如 $N(a)=0$ 对所有 $a$ 成立。因为对 $b\neq0$，总有 $a=qb+0$，其中 $q=ab^{-1}$。
	
	\item 整数环 $\ZZ$ 是 Euclidean整环，范数为 $N(a)=|a|$。这就是初等算术中的长除法。若 $a,b$ 是非零整数且先设 $b>0$，则半开区间 $[nb,(n+1)b)$，$n\in\ZZ$，划分实数轴，故 $a$ 落在某个区间中，设 $a\in[kb,(k+1)b)$。令 $q=k$，则 $r=a-qb$ 满足 $0\leq r<|b|$。若 $b<0$，对 $-b>0$ 应用刚才的论证即可。若 $a$ 不是 $b$ 的倍数，通常有两个可能的商与余数。例如
	\[
	5=2\cdot2+1=3\cdot2-1.
	\]
	若要求余数非负，则商和余数唯一。
	
	\item 若 $F$ 是域，则多项式环 $F[x]$ 是 Euclidean整环，范数为多项式的次数。其除法算法就是多项式长除法；证明与整数情形类似，将在下一章给出。这里系数必须来自域，因为除法算法最终依赖于能够除以任意非零系数。我们将在第 2 节证明，若 $R$ 不是域，则 $R[x]$ 不是 Euclidean整环。
	
	\item 第 7.1 节中的二次整数环 $\mcO$ 都是整环，并带有由域范数绝对值得到的范数；但一般地，$\mcO$ 并不关于该范数（也不关于任何范数）成为 Euclidean整环。不过 Gaussian 整数环 $\ZZ[i]$ 是 Euclidean整环，其范数为
	\[
	N(a+bi)=a^2+b^2.
	\]
	设 $\alpha=a+bi$，$\beta=c+di\neq0$ 是 $\ZZ[i]$ 中元素。在域 $\QQ(i)$ 中
	\[
	{\alpha\over\beta}=r+si,
	\qquad
	r={ac+bd\over c^2+d^2},\quad s={bc-ad\over c^2+d^2}.
	\]
	取整数 $p,q$ 分别最接近有理数 $r,s$，则 $|r-p|\leq1/2$ 且 $|s-q|\leq1/2$。令
	\[
	\theta=(r-p)+(s-q)i,
	\qquad \gamma=\beta\theta.
	\]
	于是 $\gamma=\alpha-(p+qi)\beta\in\ZZ[i]$，且
	\[
	N(\gamma)=N(\theta)N(\beta)\leq {1\over2}N(\beta)<N(\beta).
	\]
	所以
	\[
	\alpha=(p+qi)\beta+\gamma
	\]
	给出了除法算法。该算法很显式；商 $p+qi$ 可由 $r,s$ 很快确定。不过商不一定唯一。
	
	同样的证明也可用于证明当 $D=-2,-3,-7,-11$ 时，相应的二次整数环关于域范数是 Euclidean整环。稍后我们将看到 $\ZZ[\sqrt{-5}]$ 不是 Euclidean整环；而 $\ZZ[(1+\sqrt{-19})/2]$ 对任何范数都不是 Euclidean整环的证明在本节末给出。
	
	\item 离散赋值环也是 Euclidean整环。设 $K$ 是域，$v:K^\times\to\ZZ$ 是离散赋值，即满足
	\[
	v(ab)=v(a)+v(b),\qquad v \text{ 满射},\qquad
	v(x+y)\geq \min\{v(x),v(y)\}
	\]
	（当 $x+y\neq0$ 时）。集合
	\[
	\{x\in K^\times\mid v(x)\geq0\}\cup\{0\}
	\]
	是 $K$ 的子环，称为赋值环。若整环 $R$ 是某个离散赋值在其分式域上的赋值环，则称 $R$ 是离散赋值环。例如，对固定素数 $p$，所有分母与 $p$ 互素的有理数组成 $\QQ$ 中的离散赋值环。
	
	在离散赋值环中，令 $N(0)=0$，并在非零元素上令 $N=v$，则 $N$ 给出 Euclidean 算法：若 $N(a)<N(b)$，则 $a=0\cdot b+a$；若 $N(a)\geq N(b)$，则 $q=ab^{-1}\in R$，所以 $a=qb+0$。
\end{enumerate}

\medskip
\noindent\textbf{命题 1.}
Euclidean整环中的每个理想都是主理想。更确切地说，若 $I$ 是 Euclidean整环 $R$ 中的任一非零理想，则
\[
I=(d),
\]
其中 $d$ 是 $I$ 中任一具有最小范数的非零元素。

\begin{proof}
	若 $I=(0)$，无须证明。否则取 $I$ 中非零且范数最小的元素 $d$，其存在性来自 $\ZZ$ 的良序性。显然 $(d)\subseteq I$。反过来，取 $a\in I$，由除法算法可写
	\[
	a=qd+r,
	\qquad r=0\text{ 或 }N(r)<N(d).
	\]
	由于 $r=a-qd\in I$，$d$ 的范数最小性迫使 $r=0$。于是 $a=qd\in(d)$。故 $I=(d)$。
\end{proof}

命题 1 可用于证明某些整环不可能是 Euclidean整环：只要找出一个非主理想即可。

\paragraph{例}
\begin{enumerate}[label=(\arabic*)]
	\item 设 $R=\ZZ[x]$。理想 $(2,x)$ 不是主理想，所以 $\ZZ[x]$ 不是 Euclidean整环；尽管 $\QQ[x]$ 是 Euclidean整环。
	
	\item 设 $R=\ZZ[\sqrt{-5}]$，其域范数为 $N(a+b\sqrt{-5})=a^2+5b^2$。考虑理想
	\[
	I=(3,2+\sqrt{-5}).
	\]
	若 $I=(a+b\sqrt{-5})$ 为主理想，则 $3$ 与 $2+\sqrt{-5}$ 都被 $a+b\sqrt{-5}$ 整除。取范数得
	\[
	9=N(\alpha)(a^2+5b^2),
	\]
	从而 $a^2+5b^2$ 只能是 $1,3,9$。值为 $3$ 不可能；值为 $9$ 会推出生成元为 $\pm3$，与 $2+\sqrt{-5}$ 不被 $3$ 整除矛盾；值为 $1$ 会推出 $I=R$，从而 $1=3\gamma+(2+\sqrt{-5})\delta$，再乘以 $2-\sqrt{-5}$ 可推出 $2-\sqrt{-5}$ 被 $3$ 整除，也矛盾。因此 $I$ 不是主理想，$R$ 不是 Euclidean整环。
\end{enumerate}

Euclidean算法的一个基本后果是产生最大公因子。我们先在一般交换环中精确定义最大公因子。

\noindent\textbf{定义.}
设 $R$ 是交换环，$a,b\in R$ 且 $b\neq0$。
\begin{enumerate}[label=(\arabic*)]
	\item 若存在 $x\in R$ 使 $a=bx$，则称 $a$ 是 $b$ 的倍数，或称 $b$ 整除 $a$，记作 $b\mid a$。
	\item $a$ 与 $b$ 的一个\emph{最大公因子}是一个非零元素 $d$，满足：
	\begin{enumerate}[label=(\roman*)]
		\item $d\mid a$ 且 $d\mid b$；
		\item 若 $d'\mid a$ 且 $d'\mid b$，则 $d'\mid d$。
	\end{enumerate}
	最大公因子记为 $\gcd(a,b)$，或简单记作 $(a,b)$。
\end{enumerate}

注意 $b\mid a$ 当且仅当 $a\in(b)$，当且仅当 $(a)\subseteq(b)$。若 $d$ 是 $a,b$ 的公共因子，则 $(d)$ 包含 $a,b$，从而包含由 $a,b$ 生成的理想。因此，最大公因子（若存在）可以看作包含 $a,b$ 所生成理想的唯一最小主理想的一个生成元。

\noindent\textbf{命题 2.}
若 $a,b$ 是交换环 $R$ 中非零元素，且由 $a,b$ 生成的理想是主理想 $(d)$，则 $d$ 是 $a$ 与 $b$ 的最大公因子。

这解释了为什么符号 $(a,b)$ 常同时表示由 $a,b$ 生成的理想以及 $a,b$ 的最大公因子。一个整环若每个二元生成理想 $(a,b)$ 都是主理想，称为 Bezout 整环。

\noindent\textbf{命题 3.}
设 $R$ 是整环。若两个元素 $d,d'$ 生成同一个主理想，即 $(d)=(d')$，则 $d'=ud$，其中 $u$ 是 $R$ 的单位。特别地，若 $d,d'$ 都是 $a,b$ 的最大公因子，则它们相差一个单位因子。

\begin{proof}
	可假设 $d,d'$ 非零。由 $d\in(d')$ 与 $d'\in(d)$，存在 $x,y\in R$ 使
	\[
	d=xd',\qquad d'=yd.
	\]
	故 $d=xyd$，因 $R$ 是整环且 $d\neq0$，得 $xy=1$。因此 $x,y$ 都是单位。第二个断言由最大公因子的相互整除立刻推出。
\end{proof}

\noindent\textbf{定理 4.}
设 $R$ 是 Euclidean整环，$a,b$ 是 $R$ 中非零元素。令 $d=r_n$ 为对 $a,b$ 进行 Euclidean 算法时最后一个非零余数。则
\begin{enumerate}[label=(\arabic*)]
	\item $d$ 是 $a$ 与 $b$ 的最大公因子；
	\item 主理想 $(d)$ 就是由 $a,b$ 生成的理想。特别地，存在 $x,y\in R$ 使
	\[
	d=ax+by.
	\]
\end{enumerate}

\begin{proof}
	由命题 1，$(a,b)$ 是主理想。只需证明 $r_n$ 生成此理想。先从 Euclidean 算法最后一个方程向前追踪整除关系，可得 $r_n\mid r_{n-1}$，继而归纳得到 $r_n\mid b$ 与 $r_n\mid a$。另一方面，从第一个方程开始向后归纳：$r_0=a-q_0b\in(a,b)$，$r_1=b-q_1r_0\in(a,b)$，若 $r_{k-1},r_k\in(a,b)$，则
	\[
	r_{k+1}=r_{k-1}-q_{k+1}r_k\in(a,b).
	\]
	故 $r_n\in(a,b)$。于是 $(a,b)=(r_n)$，结论成立。
\end{proof}

例如，令 $a=2210$，$b=1131$。Euclidean 算法为
\begin{align*}
	2210&=1\cdot1131+1079,\\
	1131&=1\cdot1079+52,\\
	1079&=20\cdot52+39,\\
	52&=1\cdot39+13,\\
	39&=3\cdot13.
\end{align*}
所以 $\gcd(2210,1131)=13$。回代可得
\[
13=(-22)\cdot2210+43\cdot1131.
\]

整数 Euclidean 算法非常快：计算两个整数 $a,b$ 的最大公因子的步数至多为较小数位数的 $5$ 倍。换言之，该算法关于整数大小是对数级的。

由 $\gcd(a,b)=d$ 还可得到一次不定方程
\[
ax+by=N
\]
有整数解当且仅当 $d\mid N$。若 $(x_0,y_0)$ 是一个解，则全体解为
\[
x=x_0+m{b\over d},
\qquad
y=y_0-m{a\over d},
\qquad m\in\ZZ.
\]

本节最后给出另一个有时能用来证明某整环不是 Euclidean整环的判别准则。\footnote{这里以及下一节部分材料采用 J.C. Wilson, \emph{A principal ideal ring that is not a Euclidean ring}, Math. Mag. 46 (1973), pp. 34--38 对 Th. Motzkin 思想的阐述，并使用 Kenneth S. Williams, \emph{Note on non-Euclidean Principal Ideal Domains}, Math. Mag. 48 (1975), pp. 176--177 中的一个简化。}对任意整环，记
\[
\overline R=R^\times\cup\{0\}
\]
若 $u\in R-\overline R$，且对每个 $x\in R$，存在 $z\in\overline R$ 使 $u\mid x-z$，则称 $u$ 为\emph{万能边除子}。也就是说，对 $u$ 有一种弱形式的除法算法：每个 $x$ 都可写成
\[
x=qu+z,
\]
其中 $z$ 是 $0$ 或单位。

\noindent\textbf{命题 5.}
设 $R$ 是不是域的整环。若 $R$ 是 Euclidean整环，则 $R$ 中存在万能边除子。

\begin{proof}
	设 $R$ 关于范数 $N$ 是 Euclidean整环。取 $R-\overline R$ 中范数最小的元素 $u$。对任意 $x\in R$，写
	\[
	x=qu+r,
	\]
	其中 $r=0$ 或 $N(r)<N(u)$。由 $u$ 的选择，余数 $r$ 必在 $\overline R$ 中。因此 $u$ 是万能边除子。
\end{proof}

\paragraph{例}
令
\[
R=\ZZ\left[{1+\sqrt{-19}\over2}\right].
\]
本节末的命题可用来证明它关于任何范数都不是 Euclidean整环。事实上，$R$ 中单位只有 $\pm1$，所以 $\overline R=\{0,\pm1\}$。设 $u$ 是万能边除子。对 $x=2$，$u$ 必整除 $2,1$ 或 $3$ 中的某个非零元素；因此 $u$ 只能是 $\pm2$ 或 $\pm3$。但取 $x=(1+\sqrt{-19})/2$，可以直接检查 $x,x\pm1$ 都不被 $\pm2$ 或 $\pm3$ 整除。因此不存在万能边除子，$R$ 不是 Euclidean整环。

\paragraph{练习}
\begin{enumerate}[label=\arabic*.]
	\item 对下列五组整数 $a$ 与 $b$，求它们的最大公因子 $d$，并把 $d$ 写成 $a$ 与 $b$ 的线性组合 $ax+by$。
	\begin{enumerate}[label=(\alph*)]
		\item $a=20,\ b=13$。
		\item $a=69,\ b=372$。
		\item $a=11391,\ b=5673$。
		\item $a=507885,\ b=60808$。
		\item $a=91442056588823,\ b=779086434385541$（Euclidean 算法对这两个整数只需 $7$ 步）。
	\end{enumerate}
	
	\item 对下列每一对整数 $a$ 与 $n$，证明 $a$ 与 $n$ 互素，并求 $a$ 模 $n$ 的逆元（参见绪论章第 3 节）。
	\begin{enumerate}[label=(\alph*)]
		\item $a=13,\ n=20$。
		\item $a=69,\ n=89$。
		\item $a=1891,\ n=3797$。
		\item $a=6003722857,\ n=77695236973$（Euclidean 算法对这两个整数只需 $3$ 步）。
	\end{enumerate}
	
	\item 设 $R$ 是 Euclidean整环。设 $m$ 是所有非零元素范数的集合中的最小整数。证明 $R$ 中每个范数为 $m$ 的非零元素都是单位。推出：若存在范数为零的非零元素，则它必为单位。
	
	\item 设 $R$ 是 Euclidean整环。
	\begin{enumerate}[label=(\alph*)]
		\item 证明：若 $(a,b)=1$ 且 $a\mid bc$，则 $a\mid c$。更一般地，证明：若 $a\mid bc$ 且 $a,b$ 非零，则
		\[
		{a\over (a,b)}\mid c.
		\]
		\item 考虑 Diophantine 方程
		\[
		ax+by=N,
		\]
		其中 $a,b,N$ 为整数且 $a,b$ 非零。设 $x_0,y_0$ 是一个解：
		\[
		ax_0+by_0=N.
		\]
		证明此方程的全体解为
		\[
		x=x_0+m\,{b\over (a,b)},\qquad
		y=y_0-m\,{a\over (a,b)},
		\]
		其中 $m$ 遍历所有整数。\emph{提示：}若 $x,y$ 是 $ax+by=N$ 的解，证明
		\[
		a(x-x_0)=b(y_0-y)
		\]
		并使用 (a)。
	\end{enumerate}
	
	\item 求下列方程的全部整数解：
	\begin{enumerate}[label=(\alph*)]
		\item $2x+4y=5$。
		\item $17x+29y=31$。
		\item $85x+145y=505$。
	\end{enumerate}
	
	\item （邮票问题）设 $a,b$ 是两个互素正整数。证明：任意充分大的正整数 $N$ 都可以写成 $a$ 与 $b$ 的非负整数线性组合 $ax+by$；也就是说，存在整数 $N_0$，使得对所有 $N\geq N_0$，方程
	\[
	ax+by=N
	\]
	都有非负整数解 $x,y$。事实上，证明整数 $ab-a-b$ 不能写成 $a$ 与 $b$ 的正线性组合，但每个大于 $ab-a-b$ 的整数都是 $a$ 与 $b$ 的正线性组合（因此每个大于 $ab-a-b$ 的“邮资”都可以只用面值为 $a$ 与 $b$ 的邮票得到）。
	
	\item 用 Euclidean 算法求 $\ZZ[i]$ 中理想 $(85,1+13i)$ 的生成元，即求 $85$ 与 $1+13i$ 的最大公因子。对理想 $(47-13i,53+56i)$ 做同样的计算。
	
	\noindent 已知（但证明并不容易）$D=-1,-2,-3,-7,-11,-19,-43,-67,-163$ 是仅有的负整数 $D$，使得 $\mcO$ 中每个理想都是主理想（即按下一节术语，$\mcO$ 是 P.I.D.）。下一题的结果精确判定了哪些 $D<0$ 的二次整数环是 Euclidean整环。
	
	\item 设 $F=\QQ(\sqrt D)$ 是二次域，$\mcO$ 是其对应的二次整数环，$N$ 是第 7.1 节中的域范数。
	\begin{enumerate}[label=(\alph*)]
		\item 设 $D=-1,-2,-3,-7$ 或 $-11$。证明 $\mcO$ 关于 $N$ 是 Euclidean整环。\emph{提示：}修改正文中对 $\ZZ[i]$（$D=-1$）的证明。对 $D=-3,-7,-11$，证明 $F$ 中每个元素都与 $\mcO$ 中某个元素相差一个范数至多为
		\[
		{(1+|D|)^2\over 16|D|}
		\]
		的元素，而这个数对这些 $D$ 都小于 $1$。把 $\mcO$ 中的点画在 $\CC$ 中可能有帮助。
		\item 设 $D=-43,-67$ 或 $-163$。证明 $\mcO$ 关于任何范数都不是 Euclidean整环。\emph{提示：}套用正文中 $D=-19$ 的证明。
	\end{enumerate}
	
	\item 证明二次整数域 $\QQ(\sqrt2)$ 的整数环 $\mcO$ 关于第 7.1 节中域范数绝对值给出的范数是 Euclidean整环。
	
	\item 证明：对 $\ZZ[i]$ 的任意非零理想 $I$，商环 $\ZZ[i]/I$ 是有限的。\emph{提示：}利用 $I=(\alpha)$（其中 $\alpha\neq0$）以及此 Euclidean整环中的除法算法，说明 $I$ 的每个陪集都可由范数小于 $N(\alpha)$ 的元素代表。
	
	\item 设 $R$ 是含 $1$ 的交换环，$a,b$ 是 $R$ 中非零元素。$a$ 与 $b$ 的一个最小公倍数是指 $R$ 中的一个元素 $e$，满足：
	\begin{enumerate}[label=(\roman*)]
		\item $a\mid e$ 且 $b\mid e$；
		\item 若 $a\mid e'$ 且 $b\mid e'$，则 $e\mid e'$。
	\end{enumerate}
	\begin{enumerate}[label=(\alph*)]
		\item 证明：$a$ 与 $b$ 的最小公倍数（若存在）是包含在 $(a)\cap(b)$ 中的唯一最大主理想的一个生成元。
		\item 推出：Euclidean整环中任意两个非零元素都有最小公倍数，并且在乘以单位的意义下唯一。
		\item 证明：在 Euclidean整环中，$a$ 与 $b$ 的最小公倍数为
		\[
		{ab\over (a,b)},
		\]
		其中 $(a,b)$ 是 $a$ 与 $b$ 的最大公因子。
	\end{enumerate}
	
	\item （一种公钥密码）设 $N$ 为正整数。设 $M$ 是与 $N$ 互素的整数，$d$ 是与 $\varphi(N)$ 互素的整数，其中 $\varphi$ 表示 Euler 的 $\varphi$ 函数。证明：若
	\[
	M_1\equiv M^d\pmod N,
	\]
	则
	\[
	M\equiv M_1^{d'}\pmod N,
	\]
	其中 $d'$ 是 $d$ 模 $\varphi(N)$ 的逆元：
	\[
	dd'\equiv1\pmod{\varphi(N)}.
	\]
	
	\noindent\emph{说明：}这个结果是标准公钥密码的基础。设 $N=pq$ 是两个不同大素数的乘积（例如每个约有 $100$ 位数字）。若 $M$ 是一条消息，则 $M_1\equiv M^d\pmod N$ 是 $M$ 的加密版本，而通过计算 $M_1^{d'}\pmod N$ 可以解密（即使 $M,N$ 很大，这些幂也可通过连续平方相当容易地计算）。$N$ 与 $d$ 的值（但不是 $p,q$）公开，因此任何拥有消息 $M$ 的人都可以发送其编码消息 $M^d\pmod N$。要解码，似乎必须确定 $d'$，这又需要确定
	\[
	\varphi(N)=\varphi(pq)=(p-1)(q-1)
	\]
	（不过尚未有人证明不存在其他解码方案）。这种方法作为密码的成功依赖于分解 $N$ 为素数乘积的必要性，而目前没有足够有效的算法可完成这一点（例如最朴素的检查所有小于等于 $\sqrt N$ 的因子的办法，在这里需要约 $10^{100}$ 次计算；即使每秒 $100$ 亿次计算，也大约需要 $300$ 年，而且当然还可以继续增大 $p,q$ 的大小）。
\end{enumerate}


\section{主理想整环（P.I.D.s）}

\noindent\textbf{定义.}
\emph{主理想整环}（Principal Ideal Domain，P.I.D.）是每个理想都是主理想的整环。

命题 1 已经证明每个 Euclidean整环都是主理想整环。因此，所有关于主理想整环的结果自动适用于 Euclidean整环。

\paragraph{例}
\begin{enumerate}[label=(\arabic*)]
	\item 如命题 1 后面所说，整数环 $\ZZ$ 是 P.I.D.。我们在第 7.4 节看到，多项式环 $\ZZ[x]$ 含有非主理想，因此不是 P.I.D.。
	\item 命题 1 后的例 2 说明，二次整数环 $\ZZ[\sqrt{-5}]$ 不是 P.I.D.；事实上，理想 $(3,1+\sqrt{-5})$ 是非主理想。两个非主理想的乘积可以是主理想，例如
	\[
	(3,1+\sqrt{-5})(3,1-\sqrt{-5})=(3)
	\]
	（参见练习 5 以及下面命题 12 之前的例子）。
\end{enumerate}

并非每个主理想整环都是 Euclidean整环。下面将证明，上一节已经证明不是 Euclidean整环的二次整数环
\[
\ZZ\left[{1+\sqrt{-19}\over2}\right]
\]
仍然是一个 P.I.D.。

从理想论的观点看，主理想整环是在域（其理想只有平凡的 $(0)$ 和 $(1)$）之外自然要研究的一类环。Euclidean整环所享有的许多性质在主理想整环中也成立。不过 Euclidean整环相对于主理想整环的一个显著优点是：虽然在两种情形下最大公因子都存在，但在 Euclidean整环中有一个算法来计算它们。因此，依赖于最大公因子存在性的结果往往可以在更大的主理想整环类中证明；然而具体例子的计算（即这些结果的具体应用）在有 Euclidean算法可用时通常更有效。

我们收集上一节中已经证明的关于最大公因子的若干事实。

\medskip
\noindent\textbf{命题 6.}
设 $R$ 是主理想整环，$a,b$ 是 $R$ 中非零元素。令 $d$ 为由 $a,b$ 生成的主理想的一个生成元。则
\begin{enumerate}[label=(\arabic*)]
	\item $d$ 是 $a,b$ 的最大公因子；
	\item $d$ 可写成 $a,b$ 的 $R$-线性组合，也就是说，存在 $x,y\in R$ 使
	\[
	d=ax+by;
	\]
	\item $d$ 在乘以 $R$ 中单位的意义下唯一。
\end{enumerate}

\begin{proof}
	这只是命题 2 与命题 3。
\end{proof}

回忆，极大理想总是素理想，但反过来一般不真。我们在第 7.4 节观察到，$\ZZ$ 的每个非零素理想都是极大理想。如下命题说明，这个有用事实在任意主理想整环中成立。

\medskip
\noindent\textbf{命题 7.}
主理想整环中的每个非零素理想都是极大理想。

\begin{proof}
	设 $(p)$ 是主理想整环 $R$ 中的非零素理想，并设 $I=(m)$ 是包含 $(p)$ 的任一理想。我们必须证明 $I=(p)$ 或 $I=R$。由于 $p\in(m)$，有 $p=rm$ 对某个 $r\in R$ 成立。因为 $(p)$ 是素理想且 $rm\in(p)$，所以 $r$ 或 $m$ 必在 $(p)$ 中。若 $m\in(p)$，则 $(p)=(m)=I$。另一方面，若 $r\in(p)$，写 $r=ps$。此时
	\[
	p=rm=psm,
	\]
	所以 $sm=1$（回忆 $R$ 是整环），故 $m$ 是单位，于是 $I=R$。
\end{proof}

如前所述，若 $F$ 是域，则多项式环 $F[x]$ 是 Euclidean整环，从而也是主理想整环（这将在下一章证明）。其逆命题也成立。直观地说，若 $I$ 是 $R$ 中的一个理想（例如 $\ZZ$ 中的理想 $(2)$），则 $R[x]$ 中的理想 $(I,x)$（例如 $\ZZ[x]$ 中的理想 $(2,x)$）通常比 $I$ 本身多需要一个生成元，因此一般不是主理想。

\medskip
\noindent\textbf{推论 8.}
若 $R$ 是任意交换环，并且多项式环 $R[x]$ 是主理想整环（或 Euclidean整环），则 $R$ 必为域。

\begin{proof}
	假设 $R[x]$ 是主理想整环。因为 $R$ 是 $R[x]$ 的子环，所以 $R$ 必为整环（回忆 $R[x]$ 有单位当且仅当 $R$ 有单位）。理想 $(x)$ 是 $R[x]$ 中的非零素理想，因为 $R[x]/(x)$ 同构于整环 $R$。由命题 7，$(x)$ 是极大理想，于是商环 $R$ 由第 7.4 节命题 12 是域。
\end{proof}

本节最后一个结果将用来证明并非每个 P.I.D. 都是 Euclidean整环，并把主理想性质同 Euclidean 条件的另一个弱化形式联系起来。

\medskip
\noindent\textbf{定义.}
若 $N$ 是正范数，并且对每个非零 $a,b\in R$，或者 $a$ 是理想 $(b)$ 的元素，或者理想 $(a,b)$ 中有一个范数严格小于 $b$ 的范数的非零元素，则称 $N$ 是一个 \emph{Dedekind--Hasse范数}。也就是说，或者 $b$ 在 $R$ 中整除 $a$，或者存在 $s,t\in R$ 使
\[
0<N(sa-tb)<N(b).
\]

注意，如果对某个正范数 $N$，$R$ 是 Euclidean的，则总能用 $s=1$ 满足 Dedekind--Hasse 条件，所以这确实是 Euclidean 条件的一个弱化。

\medskip
\noindent\textbf{命题 9.}
整环 $R$ 是 P.I.D. 当且仅当 $R$ 有一个 Dedekind--Hasse范数。\footnote{Dedekind--Hasse 范数在 $R$ 上推出 $R$ 是 P.I.D.（当 $R$ 是代数整数环时还等价）是 Dedekind 与 Hasse 的经典判别准则，见 \emph{\"Uber eindeutige Zerlegung in Primelemente oder in Primhauptideale in Integrit\"atsbereichen}, Jour. f\"ur die Reine und Angew. Math. 159 (1928), pp. 3--12。其逆命题一般成立这一观察较新，见 John Greene, \emph{Principal Ideal Domains are almost Euclidean}, Amer. Math. Monthly 104 (1997), pp. 154--156。}

\begin{proof}
	令 $I$ 是 $R$ 的任一非零理想，取 $I$ 中满足 $N(b)$ 最小的非零元素 $b$。设 $a$ 是 $I$ 中任一非零元素，则理想 $(a,b)$ 包含于 $I$。于是 Dedekind--Hasse 条件以及 $b$ 的最小性说明 $a\in(b)$，所以 $I=(b)$ 是主理想。反方向将在下一节证明（推论 16）。
\end{proof}

\paragraph{例}
令
\[
R=\ZZ\left[{1+\sqrt{-19}\over2}\right]
\]
为上一节末所考虑的二次整数环。我们说明定义在 $R$ 上的正的域范数
\[
N\left(a+b{1+\sqrt{-19}\over2}\right)=a^2+ab+5b^2
\]
是 Dedekind--Hasse范数；于是由命题 9 与上一节结果可知，$R$ 是 P.I.D. 但不是 Euclidean整环。

设 $\alpha,\beta$ 是 $R$ 的非零元素且 $\alpha/\beta\notin R$。我们必须证明存在 $s,t\in R$ 使
\[
0<N(s\alpha-t\beta)<N(\beta),
\]
由域范数的乘法性，这等价于
\[
0<N\left({\alpha\over\beta}s-t\right)<1. \tag{*}
\]
写
\[
{\alpha\over\beta}={a+b\sqrt{-19}\over c}\in \QQ[\sqrt{-19}],
\]
其中整数 $a,b,c$ 没有公因子且 $c>1$（因为假设 $\beta$ 不整除 $\alpha$）。既然 $a,b,c$ 没有公因子，存在整数 $x,y,z$ 使 $ax+by+cz=1$。把
\[
ay-19bx=cq+r
\]
写成商为 $q$、余数为 $r$ 且 $|r|\leq c/2$ 的形式，并令
\[
s=y+x\sqrt{-19},\qquad t=q-z\sqrt{-19}.
\]
则直接计算给出
\[
0<N\left({\alpha\over\beta}s-t\right)
={ (ay-19bx-cq)^2+19(ax+by+cz)^2\over c^2}
\leq {1\over4}+{19\over c^2},
\]
因此当 $c\geq5$ 时，上述 $s,t$ 满足 $(*)$。

若 $c=2$，则 $a,b$ 中一个为偶数而另一个为奇数（否则 $\alpha/\beta\in R$），这时直接检查可知
\[
s=1,
\qquad
 t={ (a-1)+b\sqrt{-19}\over2}
\]
是 $R$ 的元素并满足 $(*)$。

若 $c=3$，则整数 $a^2+19b^2$ 不能被 $3$ 整除（模 $3$ 后这是 $a^2+b^2$，容易看出它模 $3$ 为 $0$ 当且仅当 $a,b$ 都模 $3$ 为 $0$；但那样 $a,b,c$ 有公因子）。写
\[
a^2+19b^2=3q+r,
\qquad r=1\text{ 或 }2.
\]
于是再次直接检查可知
\[
s=a-b\sqrt{-19},
\qquad t=q
\]
是 $R$ 的元素并满足 $(*)$。

最后，若 $c=4$，则 $a,b$ 不全为偶数。若 $a,b$ 中一个为偶数而另一个为奇数，则 $a^2+19b^2$ 是奇数，所以可写
\[
a^2+19b^2=4q+r,
\qquad 0<r<4,
\]
其中 $q,r\in\ZZ$。这时 $s=a-b\sqrt{-19}$ 与 $t=q$ 满足 $(*)$。若 $a,b$ 都为奇数，则
\[
a^2+19b^2\equiv 1+3 \pmod 8,
\]
所以可写 $a^2+19b^2=8q+4$，其中 $q\in\ZZ$。这时
\[
s={a-b\sqrt{-19}\over2},
\qquad t=q
\]
是 $R$ 中元素并满足 $(*)$。

\paragraph{练习}
\begin{enumerate}[label=\arabic*.]
	\item 证明：在主理想整环中，两个理想 $(a)$ 与 $(b)$ 互素（参见第 7.6 节）当且仅当 $a$ 与 $b$ 的一个最大公因子是 $1$（此时称 $a$ 与 $b$ \emph{互素}或\emph{相对素}）。
	
	\item 证明：P.I.D. 中任意两个非零元素都有最小公倍数（参见第 1 节练习 11）。
	
	\item 证明：P.I.D. 对素理想取商仍是 P.I.D.。
	
	\item 设 $R$ 是整环。证明：若下列两个条件成立，则 $R$ 是主理想整环：
	\begin{enumerate}[label=(\roman*)]
		\item $R$ 中任意两个非零元素 $a,b$ 都有最大公因子，且该最大公因子可写成 $ra+sb$ 的形式，其中 $r,s\in R$；
		\item 若 $a_1,a_2,a_3,\ldots$ 是 $R$ 中非零元素，且对所有 $i$ 都有 $a_{i+1}\mid a_i$，则存在正整数 $N$，使得对所有 $n\geq N$，$a_n$ 都是 $a_N$ 乘以某个单位。
	\end{enumerate}
	
	\item 令 $R$ 为二次整数环 $\ZZ[\sqrt{-5}]$。定义理想
	\[
	I_2=(2,1+\sqrt{-5}),\qquad
	I_3=(3,2+\sqrt{-5}),\qquad
	I_3'=(3,2-\sqrt{-5}).
	\]
	\begin{enumerate}[label=(\alph*)]
		\item 证明 $I_2,I_3,I_3'$ 都是 $R$ 中的非主理想。\emph{提示：}命题 1 后面的例 2 已对 $I_3$ 给出证明。
		\item 通过证明 $I_2^2$ 是由 $2$ 生成的主理想，即
		\[
		I_2^2=(2),
		\]
		说明两个非主理想的乘积可以是主理想。
		\item 类似地证明
		\[
		I_2I_3=(1-\sqrt{-5}),\qquad I_2I_3'=(1+\sqrt{-5})
		\]
		都是主理想。推出主理想 $(6)$ 是四个理想的乘积：
		\[
		(6)=I_2^2I_3I_3'.
		\]
	\end{enumerate}
	
	\item 设 $R$ 是整环，并假设 $R$ 中每个素理想都是主理想。本题证明 $R$ 中每个理想都是主理想，即 $R$ 是 P.I.D.。
	\begin{enumerate}[label=(\alph*)]
		\item 假设 $R$ 中非主理想的集合非空，证明这个集合在包含关系下有一个极大元（按假设它不是素理想）。\emph{提示：}使用 Zorn 引理。
		\item 令 $I$ 是一个关于“非主”性质极大的理想，且取 $a,b\in R$ 满足 $ab\in I$，但 $a\notin I$ 且 $b\notin I$。令 $I_a=(I,a)$ 为由 $I$ 与 $a$ 生成的理想，令 $I_b=(I,b)$ 为由 $I$ 与 $b$ 生成的理想，并定义
		\[
		J=\{r\in R\mid rI_a\subseteq I\}.
		\]
		证明 $I_a=(\alpha)$ 与 $J=(\beta)$ 都是 $R$ 中的主理想，并且
		\[
		I\subseteq I_b\subseteq J,
		\qquad I_aJ=(\alpha\beta)\subseteq I.
		\]
		\item 若 $x\in I$，证明 $x=s\alpha$ 对某个 $s\in J$ 成立。推出 $I=I_aJ$ 是主理想，得到矛盾；从而推出 $R$ 是 P.I.D.。
	\end{enumerate}
	
	\item 一个整环 $R$ 若满足每个由两个元素生成的理想都是主理想（即对每个 $a,b\in R$，都有 $(a,b)=(d)$ 对某个 $d\in R$ 成立），则称为 \emph{Bezout整环}（也参见第 3 节练习 11）。
	\begin{enumerate}[label=(\alph*)]
		\item 证明：整环 $R$ 是 Bezout整环当且仅当 $R$ 中每一对元素 $a,b$ 都有一个最大公因子 $d$，并且 $d$ 可写成 $a,b$ 的 $R$-线性组合：
		\[
		d=ax+by
		\]
		其中 $x,y\in R$。
		\item 证明：Bezout整环的每个有限生成理想都是主理想。\emph{提示：}在第 9.2 节与第 9.3 节的练习中可以看到不是每个理想都主的 Bezout整环。
		\item 设 $F$ 是 Bezout整环 $R$ 的分式域。证明 $F$ 的每个元素都可以写成 $a/b$ 的形式，其中 $a,b\in R$ 且 $a,b$ 互素（参见练习 1）。
	\end{enumerate}
	
	\item 证明：若 $R$ 是主理想整环，$D$ 是 $R$ 的一个乘法封闭子集，则 $D^{-1}R$ 也是 P.I.D.（参见第 7.5 节）。
\end{enumerate}


\section{唯一分解整环（U.F.D.s）}

在整数环 $\ZZ$ 的情形中，还有另一种从初等算术中熟悉的方法来确定两个元素 $a,b$ 的最大公因子，即把 $a,b$ “分解为素数”，再由此容易地确定最大公因子。这也可以推广到一类更大的环，称为唯一分解整环（U.F.D.）——我们马上给出定义。随后我们将证明
\[
\text{每个主理想整环都是唯一分解整环}
\]
从而关于唯一分解整环的每个结果都自动适用于 Euclidean整环和主理想整环。

我们先引入一些术语。

\medskip
\noindent\textbf{定义.}
设 $R$ 是整环。
\begin{enumerate}[label=(\arabic*)]
	\item 设 $r\in R$ 非零且不是单位。若每当 $r=ab$（其中 $a,b\in R$）时，$a,b$ 中至少一个必须是 $R$ 中的单位，则称 $r$ 在 $R$ 中\emph{不可约}。否则称 $r$ \emph{可约}。
	\item 非零元素 $p\in R$ 称为 $R$ 中的\emph{素元}，若由 $p$ 生成的理想 $(p)$ 是素理想。换言之，非零元素 $p$ 是素元，当且仅当 $p$ 不是单位，并且对任意 $a,b\in R$，若 $p\mid ab$，则 $p\mid a$ 或 $p\mid b$。
	\item 若 $R$ 中两个元素 $a,b$ 相差一个单位因子，即 $a=ub$ 对某个单位 $u\in R$ 成立，则称 $a,b$ 在 $R$ 中\emph{相伴}。
\end{enumerate}

\medskip
\noindent\textbf{命题 10.}
整环中的素元总是不可约元。

\begin{proof}
	设 $(p)$ 是非零素理想且 $p=ab$。于是 $ab=p\in(p)$，故由素理想的定义，$a$ 或 $b$ 中有一个（设为 $a$）属于 $(p)$。从而 $a=pr$ 对某个 $r\in R$ 成立。这推出
	\[
	p=ab=prb,
	\]
	所以 $rb=1$，故 $b$ 是单位。这说明 $p$ 不可约。
\end{proof}

一般地，不可约元不一定是素元。例如，考虑二次整数环 $R=\ZZ[\sqrt{-5}]$ 中的元素 $3$。第 1 节中的计算说明 $3$ 在 $R$ 中不可约，但 $3$ 不是素元，因为
\[
(2+\sqrt{-5})(2-\sqrt{-5})=3^2
\]
被 $3$ 整除，而 $2+\sqrt{-5}$ 与 $2-\sqrt{-5}$ 中没有一个在 $R$ 中被 $3$ 整除。

然而，若 $R$ 是主理想整环，则素元与不可约元两个概念相同。特别地，在 $\ZZ$ 与 $F[x]$（其中 $F$ 是域）中二者相同。

\medskip
\noindent\textbf{命题 11.}
在主理想整环中，非零元素是素元当且仅当它不可约。

\begin{proof}
	我们已经证明素元蕴含不可约。反过来，设 $p$ 不可约。我们必须证明 $p$ 是素元，即理想 $(p)$ 是素理想。若 $M$ 是任一包含 $(p)$ 的理想，则由假设 $M=(m)$ 是主理想。由于 $p\in(m)$，有 $p=rm$ 对某个 $r$ 成立。但 $p$ 不可约，所以按定义 $r$ 或 $m$ 是单位。这分别意味着 $(p)=(m)$ 或 $(m)=(1)$。因此包含 $(p)$ 的理想只有 $(p)$ 和 $(1)$，即 $(p)$ 是极大理想。极大理想都是素理想，证明完成。
\end{proof}

\paragraph{例}
命题 11 给出二次整数环 $\ZZ[\sqrt{-5}]$ 不是 P.I.D. 的另一个证明：因为 $3$ 在此环中不可约但不是素元。

整数环 $\ZZ$ 中的不可约元就是初等算术中熟悉的素数（以及它们的负数），而两个整数 $a,b$ 相伴当且仅当 $a=\pm b$。

在整数环 $\ZZ$ 中，任一整数 $n$ 都可写成素数（不一定互异）的乘积。理由如下。若 $n$ 本身不是素数，则由定义可以写成 $n=n_1n_2$，其中 $n_1,n_2$ 都不是单位，即都不是 $\pm1$。$n_1,n_2$ 的绝对值都小于 $n$ 本身。若二者都是素数，则已经把 $n$ 写成素数乘积。若 $n_1$ 或 $n_2$ 中有一个不是素数，则它又可分解成两个（更小的）整数。由于整数的绝对值不能无限下降，我们最终必然只剩下素数因子，于是得到 $n$ 的素数乘积分解。

例如，若 $n=2210$，上述算法如下进行：$n$ 不是素数，因为
\[
n=2\cdot1105.
\]
整数 $2$ 是素数，但 $1105$ 不是：
\[
1105=5\cdot221.
\]
整数 $5$ 是素数，但 $221$ 不是：
\[
221=13\cdot17.
\]
到此算法终止，因为 $13$ 与 $17$ 都是素数。这给出 $2210$ 的素因子分解
\[
2210=2\cdot5\cdot13\cdot17.
\]
类似地，我们得到
\[
1131=3\cdot13\cdot29.
\]
在这些例子中每个素数只出现一次，但一般当然不必如此。

在环 $\ZZ$ 中，不但每个整数 $n$ 都能写成素数乘积，而且这种分解事实上是唯一的：同一个正整数 $n$ 的任意两个素因子分解只差正素因子的书写顺序。这里限制为正整数，是为了避免把 $15$ 的分解 $(3)(5)$ 与 $(-3)(-5)$ 看作本质不同。$\ZZ$ 的这个唯一分解性质（我们很快将证明）对于整数算术非常有用。具有类似性质的一般环有一个名称。

\medskip
\noindent\textbf{定义.}
\emph{唯一分解整环}（Unique Factorization Domain，U.F.D.）是一个整环 $R$，其中每个非零非单位元素 $r\in R$ 都有如下两个性质：
\begin{enumerate}[label=(\roman*)]
	\item $r$ 可写成 $R$ 中有限个不可约元 $p_i$ 的乘积（不一定互异）：
	\[
	r=p_1p_2\cdots p_n;
	\]
	\item (i) 中的分解在相伴意义下唯一：即若
	\[
	r=q_1q_2\cdots q_m
	\]
	是 $r$ 的另一个不可约元分解，则 $m=n$，并且在重新编号之后，对 $i=1,2,\ldots,n$，$p_i$ 与 $q_i$ 相伴。
\end{enumerate}

\paragraph{例}
\begin{enumerate}[label=(\arabic*)]
	\item 域 $F$ 平凡地是唯一分解整环，因为每个非零元素都是单位，所以没有元素需要验证性质 (i) 和 (ii)。
	\item 如上所述，我们很快将证明每个主理想整环都是唯一分解整环（因此特别地，$\ZZ$ 与 $F[x]$，其中 $F$ 是域，都是唯一分解整环）。
	\item 下一章还将证明：只要 $R$ 本身是唯一分解整环，多项式环 $R[x]$ 也是唯一分解整环（这与主理想整环或 Euclidean整环的性质形成对比，后两种性质不会从 $R$ 传递到 $R[x]$）。这个结果连同上一例将说明 $\ZZ[x]$ 是唯一分解整环。
	\item Gaussian 整数的子环
	\[
	R=\ZZ[2i]=\{a+2bi\mid a,b\in\ZZ\},\qquad i^2=-1,
	\]
	是整环但不是唯一分解整环（这类环在第 7.1 节练习 23 中已经引入）。元素 $2$ 与 $2i$ 是不可约元，并且在 $R$ 中不相伴，因为 $i\notin R$；而
	\[
	4=2\cdot2=(-2i)(2i)
	\]
	给出 $R$ 中的两个不同分解。也可直接检查 $2i$ 在 $R$ 中不可约但不是素元（因为 $R/(2i)\cong\ZZ/4\ZZ$）。在更大的 Gaussian 整数环 $\ZZ[i]$（它是唯一分解整环）中，$2$ 与 $2i$ 相伴，因为 $i$ 在这个较大环中是单位。我们将在第 9.3 节末给出 $\ZZ[2i]$ 不是唯一分解整环的稍有不同的证明（该证明无需检查 $2$ 与 $2i$ 不可约）。
	\item 二次整数环 $\ZZ[\sqrt{-5}]$ 是另一个不是唯一分解整环的整环，因为
	\[
	6=2\cdot3=(1+\sqrt{-5})(1-\sqrt{-5})
	\]
	给出 $6$ 分解为不可约元的两个不同分解。$\ZZ[\sqrt{-5}]$ 中的主理想 $(6)$ 可以写成四个非主素理想的乘积：
	\[
	(6)=P_2^2P_3P_3',
	\]
	而元素 $6$ 的两个不同分解可解释为该理想乘积重新组合为主理想乘积的两种方式：$P_2^2=(2)$ 与 $P_3P_3'=(3)$ 的乘积，以及 $P_2P_3=(1+\sqrt{-5})$ 与 $P_2P_3'=(1-\sqrt{-5})$ 的乘积（参见练习 8）。
	
	虽然二次整数环 $\mcO$ 中的元素未必有唯一分解，但有一个定理（推论 16.16）断言 $\mcO$ 中每个理想都能唯一地写成素理想的乘积。理想唯一分解为素理想乘积的性质，对于代数数域的整数环（其中二次整数环就是例子）一般成立，并导向第 16 章考虑的 Dedekind整环概念。历史上，正是数论中代数整数环里的元素分解为不可约元的唯一性失败，最初导致了理想的定义。由此产生的这些环中分解为素理想的唯一性，使理想的元素具有一种“理想的”（在“完美”或“合意”的意义上）行为，这也正是这些如今基本的代数对象术语选择的基础。
\end{enumerate}

唯一分解整环中不可约元的第一个性质是它们也都是素元。也许会认为可以由下面命题以及前面提到的定理（我们马上要证明）“每个主理想整环都是唯一分解整环”推出命题 11；然而命题 11 将用于后一定理的证明中。

\medskip
\noindent\textbf{命题 12.}
在唯一分解整环中，非零元素是素元当且仅当它不可约。

\begin{proof}
	设 $R$ 是唯一分解整环。由命题 10，素元不可约；因此只需证明每个不可约元都是素元。设 $p$ 是 $R$ 中不可约元，并假设 $p\mid ab$，其中 $a,b\in R$；我们必须证明 $p$ 整除 $a$ 或 $b$。说 $p$ 整除 $ab$，就是说 $ab=pc$ 对某个 $c\in R$ 成立。把 $a,b$ 写成不可约元乘积，由这个等式以及 $ab$ 分解为不可约元的唯一性可知，不可约元 $p$ 必与 $a$ 的分解或 $b$ 的分解中出现的某个不可约元相伴。不妨设 $p$ 与 $a$ 的分解中某个不可约元相伴，也就是说，$a$ 可写成
	\[
	a=(up)p_2\cdots p_n,
	\]
	其中 $u$ 是单位，$p_2,\ldots,p_n$ 是某个（可能为空的）不可约元集合。那么 $p\mid a$，因为
	\[
	a=pd,
	\qquad d=up_2\cdots p_n.
	\]
	证明完成。
\end{proof}

在唯一分解整环中，我们现在将互换使用“素元”和“不可约元”这两个术语；不过通常仍说 $\ZZ$ 中的“素数”和 $F[x]$ 中的“不可约多项式”。

我们将用前一命题证明：在唯一分解整环中任意两个非零元素 $a,b$ 都有最大公因子。

\medskip
\noindent\textbf{命题 13.}
设 $a,b$ 是唯一分解整环 $R$ 中两个非零元素，并设
\[
a=u p_1^{e_1}p_2^{e_2}\cdots p_n^{e_n},
\qquad
b=v p_1^{f_1}p_2^{f_2}\cdots p_n^{f_n}
\]
是 $a,b$ 的素因子分解，其中 $u,v$ 是单位，素元 $p_1,p_2,\ldots,p_n$ 互不相同，并且指数 $e_i,f_i\geq0$。则元素
\[
d=p_1^{\min(e_1,f_1)}p_2^{\min(e_2,f_2)}\cdots p_n^{\min(e_n,f_n)}
\]
（若所有指数都是 $0$，则 $d=1$）是 $a,b$ 的最大公因子。

\begin{proof}
	由于 $d$ 中出现的每个素元的指数不大于 $a$ 与 $b$ 的分解中相应指数，$d$ 整除 $a$ 且整除 $b$。为证明 $d$ 是最大公因子，令 $c$ 是 $a,b$ 的任意公共因子，并写 $c=q_1^{g_1}q_2^{g_2}\cdots q_m^{g_m}$ 为 $c$ 的素因子分解。由于每个 $q_i$ 整除 $c$，因而整除 $a$ 与 $b$，由前一命题可知 $q_i$ 必整除某个 $p_j$。特别地，在相伴意义下（即至多相差一个单位），$c$ 中出现的素元必须是 $a,b$ 中出现的素元的子集：
	\[
	\{q_1,q_2,\ldots,q_m\}\subseteq\{p_1,p_2,\ldots,p_n\}.
	\]
	类似地，$c$ 中各素元的指数也不大于 $d$ 中相应指数。因此 $c\mid d$，证明完成。
\end{proof}

\paragraph{例}
在上面的例子中，$a=2210$ 且 $b=1131$，从它们的素因子分解立即得到 $(a,b)=13$。注意，如果 $a,b$ 的素因子分解已知，则上面的命题能立刻给出它们的最大公因子；但找到这些素因子分解在计算上非常耗时。Euclidean算法是确定两个整数最大公因子的最快方法，但遗憾的是，它几乎不给出这些整数的素因子分解信息。

现在我们来到本章所介绍的几类环之间的一个主要联系。

\medskip
\noindent\textbf{定理 14.}
每个主理想整环都是唯一分解整环。特别地，每个 Euclidean整环都是唯一分解整环。

\begin{proof}
	注意第二个断言由第一个断言推出，因为 Euclidean整环是主理想整环。为证明第一个断言，设 $R$ 是主理想整环，$r$ 是 $R$ 中非零且不是单位的元素。我们必须先证明 $r$ 可写成 $R$ 中有限个不可约元的乘积，然后验证该分解在相差单位的意义下唯一。
	
	第一部分证明的方法与确定整数的素因子分解完全类似。假设 $r$ 非零且不是单位。若 $r$ 本身不可约，则完成。若不然，由定义可将 $r$ 写成乘积 $r=r_1r_2$，其中 $r_1,r_2$ 都不是单位。若这两个元素都不可约，则同样完成，因为已经把 $r$ 写成不可约元乘积。否则，至少有一个元素，设为 $r_1$，是可约的，于是可写成两个非单位元素的乘积 $r_1=r_{11}r_{12}$，如此继续。我们必须证明这个过程终止，即必然到达某一步，使得作为 $r$ 的因子得到的所有元素都是不可约元。假设不是这样。由分解 $r=r_1r_2$ 得到理想的真包含
	\[
	(r)\subset (r_1)\subset R.
	\]
	第一个包含是真包含，因为 $r_2$ 不是单位；最后一个包含是真包含，因为 $r_1$ 不是单位。由 $r_1$ 的分解类似得到
	\[
	(r)\subset (r_1)\subset (r_{11})\subset R.
	\]
	若这个分解过程在有限步后不终止，则会得到一个无限上升理想链
	\[
	(r)\subset (r_1)\subset (r_{11})\subset\cdots\subset R,
	\]
	其中所有包含都是真包含，并且选择公理保证存在这样一条无限链（参见附录 I）。
	
	现在证明主理想整环中的任何上升理想链
	\[
	I_1\subseteq I_2\subseteq\cdots\subseteq R
	\]
	最终稳定，即存在正整数 $n$，使得对所有 $k\geq n$，$I_k=I_n$。\footnote{同一个论证可用来证明更一般的陈述：在每个理想都是有限生成的交换环中，上升理想链最终稳定。第 12 章将需要这个结果，并会重复其细节。}特别地，不可能存在所有包含都是真包含的无限上升理想链。令
	\[
	I=\bigcup_{i=1}^{\infty} I_i.
	\]
	容易验证（如第 7.4 节命题 11 的证明中一样）$I$ 是理想。由于 $R$ 是主理想整环，$I$ 主生成，设 $I=(a)$。因为 $I$ 是上述理想的并，$a$ 必属于链中某个理想，设 $a\in I_n$。于是
	\[
	I_n\subseteq I=(a)\subseteq I_n,
	\]
	所以 $I=I_n$，链从 $I_n$ 起稳定。这证明 $R$ 中每个非零非单位元素都有某个不可约元分解。
	
	还需证明上述分解本质唯一。我们对元素 $r$ 的某个分解中不可约因子的个数 $n$ 归纳。若 $n=0$，则 $r$ 是单位。若有 $r=qc$（某个其他分解），其中 $q$ 不可约，则 $q$ 将整除一个单位，因而 $q$ 本身是单位，矛盾。现设 $n\geq1$，且对 $r$ 有两个乘积
	\[
	r=p_1p_2\cdots p_n=q_1q_2\cdots q_m,
	\]
	其中 $p_i,q_j$ 是（不一定互异的）不可约元。于是 $p_1$ 整除右边乘积，由命题 11，$p_1$ 必整除其中一个因子。必要时重新编号，可设 $p_1\mid q_1$。于是 $q_1=p_1u$ 对某个 $u\in R$ 成立，而由于 $q_1$ 不可约，$u$ 实际上必须是单位。因此 $p_1$ 与 $q_1$ 相伴。在整环中可消去 $p_1$（这是合法的），得到
	\[
	p_2\cdots p_n=q_2' q_3\cdots q_m,
	\]
	其中 $q_2'=uq_2$ 仍不可约（与 $q_2$ 相伴）。由关于 $n$ 的归纳假设，左边每个因子与最右边的因子逐一匹配（在相伴意义下），从而也与中间原来的因子匹配（它们只差相伴）。由于 $p_1$ 与 $q_1$（在最初重新编号后）已经证明相伴，归纳步骤和定理证明完成。
\end{proof}

\medskip
\noindent\textbf{推论 15（算术基本定理）.}
整数环 $\ZZ$ 是唯一分解整环。

\begin{proof}
	整数环 $\ZZ$ 是 Euclidean整环，因此由定理是唯一分解整环。
\end{proof}

现在可以完成命题 9 中关于整环 $R$ 是否为 P.I.D. 与是否存在 Dedekind--Hasse范数之间的等价性。

\medskip
\noindent\textbf{推论 16.}
设 $R$ 是 P.I.D.，则 $R$ 上存在一个乘法性的 Dedekind--Hasse范数。

\begin{proof}
	若 $R$ 是 P.I.D.，则 $R$ 是 U.F.D.。定义范数 $N$ 如下：令 $N(0)=0$；若 $u$ 是单位，令 $N(u)=1$；若
	\[
	a=p_1p_2\cdots p_n
	\]
	是 $a$ 的不可约元分解，则令 $N(a)=2^n$。这一定义良好，因为 $a$ 的不可约因子的个数唯一。显然 $N(ab)=N(a)N(b)$，所以 $N$ 是正的且乘法性的。为证明 $N$ 是 Dedekind--Hasse范数，设 $a,b$ 是 $R$ 中非零元素。由假设，由 $a,b$ 生成的理想是主理想，设 $(a,b)=(r)$。若 $a$ 不属于理想 $(b)$，则 $r$ 也不属于 $(b)$，即 $r$ 不被 $b$ 整除。因为 $b=xr$ 对某个 $x\in R$ 成立，所以 $x$ 不是单位，于是
	\[
	N(b)=N(x)N(r)>N(r).
	\]
	因此 $(a,b)$ 中含有一个范数严格小于 $b$ 的范数的非零元素，证明完成。
\end{proof}

\subsection*{Gaussian 整数中的分解}

我们以描述 Gaussian 整数 $\ZZ[i]$ 中的不可约元并给出它对初等数论中 Fermat 一个著名定理的应用，结束对唯一分解整环的讨论。这是特别合适的，因为对 $\ZZ[i]$ 的经典研究开启了环的代数研究。

一般地，设 $\mcO$ 是二次整数环，$N$ 是第 7.1 节中引入的相应域范数。设 $\alpha\in\mcO$ 是一个范数为 $\ZZ$ 中素数 $p$ 的元素。若 $\alpha=\beta\gamma$，其中 $\beta,\gamma\in\mcO$，则
\[
p=N(\alpha)=N(\beta)N(\gamma),
\]
所以 $N(\beta)$ 或 $N(\gamma)$ 中有一个是 $\pm1$，另一个是 $\pm p$。因为我们已经看到，$\mcO$ 中元素的范数为 $\pm1$ 当且仅当它是 $\mcO$ 中的单位，所以 $\alpha$ 的一个因子是单位。因此
\[
\text{若 }N(\alpha)\text{ 是 }\ZZ\text{ 中的 }\pm\text{素数，则 }\alpha\text{ 在 }\mcO\text{ 中不可约。}
\]

设 $\pi$ 是 $\mcO$ 中的素元，并设 $(\pi)$ 是 $\mcO$ 中由 $\pi$ 生成的理想。由于 $(\pi)$ 是 $\mcO$ 中素理想，容易检查 $(\pi)\cap\ZZ$ 是 $\ZZ$ 中素理想（若整数 $a,b$ 满足 $ab\in(\pi)$，则 $a$ 或 $b$ 是 $(\pi)$ 的元素，所以 $a$ 或 $b$ 属于 $(\pi)\cap\ZZ$）。由于 $N(\pi)$ 是 $(\pi)$ 中的非零整数，我们有
\[
(\pi)\cap\ZZ=p\ZZ
\]
对某个整数素数 $p$ 成立。由 $p\in(\pi)$ 可知，$\pi$ 在 $\mcO$ 中整除整数素数 $p$；因此，$\mcO$ 中的素元可以通过确定 $\ZZ$ 中素数在较大环 $\mcO$ 中如何分解来找到。设 $\pi$ 整除 $\mcO$ 中的素数 $p$，比如 $p=\pi\pi'$。则
\[
N(\pi)N(\pi')=N(p)=p^2.
\]
由于 $\pi$ 不是单位，只有两种可能：或者 $N(\pi)=\pm p^2$，或者 $N(\pi)=\pm p$。在前一情形，$N(\pi')=\pm1$，故 $\pi'$ 是单位，而 $p=\pi$（至多相差相伴）在 $\ZZ[i]$ 中不可约。后一情形中，$N(\pi)=N(\pi')=\pm p$，故 $\pi'$ 也不可约，并且 $p=\pi\pi'$ 是恰好两个不可约元的乘积。

现在考虑特殊情形 $D=-1$，即 Gaussian 整数 $\ZZ[i]$。我们已经看到 $\ZZ[i]$ 的单位是 $\pm1$ 与 $\pm i$。第 1 节证明了 $\ZZ[i]$ 是 Euclidean整环，所以它也是主理想整环与唯一分解整环；因此其中不可约元与素元相同，并可通过观察 $\ZZ$ 中素数在较大环 $\ZZ[i]$ 中如何分解来确定。

在此情形，若 $\alpha=a+bi$，则
\[
N(\alpha)=\alpha\overline\alpha=a^2+b^2,
\]
其中 $\overline\alpha=a-bi$ 是 $\alpha$ 的复共轭。由上面的讨论可知，素数 $p$ 在 $\ZZ[i]$ 中恰分解为两个不可约元，当且仅当
\[
p=a^2+b^2
\]
是两个整数平方之和（否则 $p$ 在 $\ZZ[i]$ 中仍不可约）。若 $p=a^2+b^2$，则 $\ZZ[i]$ 中相应的不可约元是 $a\pm bi$。

显然
\[
2=1^2+1^2
\]
是两个平方和，给出分解
\[
2=(1+i)(1-i)=-i(1+i)^2.
\]
不可约元 $1+i$ 与 $1-i=-i(1+i)$ 相伴，并且容易检查这是共轭不可约元 $a+bi$ 与 $a-bi$ 能相伴的唯一情形。

由于任何整数的平方模 $4$ 只可能同余于 $0$ 或 $1$，$\ZZ$ 中一个奇素数若是两个平方和，则必同余于 $1$ 模 $4$。因此，若 $p$ 是 $\ZZ$ 中满足 $p\equiv3\pmod4$ 的素数，则 $p$ 不是两个平方和，并且 $p$ 在 $\ZZ[i]$ 中仍不可约。

现在设 $p$ 是 $\ZZ$ 中满足 $p\equiv1\pmod4$ 的素数。我们将证明 $p$ 在 $\ZZ[i]$ 中不可能不可约；这将说明 $p=(a+bi)(a-bi)$ 在 $\ZZ[i]$ 中分解为两个不同不可约元的乘积，等价地，说明 $p=a^2+b^2$ 是两个平方和。我们先证明如下初等数论结果。

\medskip
\noindent\textbf{引理 17.}
素数 $p\in\ZZ$ 整除某个形如 $n^2+1$ 的整数，当且仅当 $p$ 是 $2$，或 $p$ 是同余于 $1$ 模 $4$ 的奇素数。

\begin{proof}
	$p=2$ 的断言是平凡的，因为 $2\mid 1^2+1$。若 $p$ 是奇素数，注意 $p\mid n^2+1$ 等价于在 $\ZZ/p\ZZ$ 中有 $n^2=-1$。这又等价于 $n$ 的剩余类在乘法群 $(\ZZ/p\ZZ)^\times$ 中有阶 $4$。因此，$p$ 整除某个形如 $n^2+1$ 的整数，当且仅当 $(\ZZ/p\ZZ)^\times$ 中含有阶为 $4$ 的元素。由 Lagrange 定理，若 $(\ZZ/p\ZZ)^\times$ 含有阶为 $4$ 的元素，则
	\[
	|(\ZZ/p\ZZ)^\times|=p-1
	\]
	被 $4$ 整除，即 $p\equiv1\pmod4$。
	
	反过来，设 $p-1$ 被 $4$ 整除。我们先说明 $(\ZZ/p\ZZ)^\times$ 中恰有一个阶为 $2$ 的元素。若 $m^2\equiv1\pmod p$，则 $p$ 整除
	\[
	m^2-1=(m-1)(m+1).
	\]
	因此 $p$ 整除 $m-1$（即 $m\equiv1\pmod p$），或者 $p$ 整除 $m+1$（即 $m\equiv-1\pmod p$），所以 $-1$ 是 $(\ZZ/p\ZZ)^\times$ 中唯一的阶为 $2$ 的剩余类。现在 Abel 群 $(\ZZ/p\ZZ)^\times$ 含有一个阶为 $4$ 的子群 $H$（例如，对子群 $\{\pm1\}$ 取商后含有阶为 $2$ 的子群，其原像就是 $(\ZZ/p\ZZ)^\times$ 中一个阶为 $4$ 的子群）。由于 Klein 四元群有三个阶为 $2$ 的元素，而 $(\ZZ/p\ZZ)^\times$——因此 $H$——只有一个阶为 $2$ 的元素，$H$ 必为四阶循环群。因此 $(\ZZ/p\ZZ)^\times$ 含有阶为 $4$ 的元素，即 $H$ 的一个生成元。
\end{proof}

\noindent\emph{说明：}稍后我们将证明（第 9.5 节推论 19）$(\ZZ/p\ZZ)^\times$ 是循环群；由此可立即看出存在阶为 $4$ 的元素当且仅当 $p-1$ 被 $4$ 整除。

由引理 17，若 $p\equiv1\pmod4$ 是素数，则存在 $n\in\ZZ$ 使得 $p$ 在 $\ZZ$ 中整除 $n^2+1$，因此当然在 $\ZZ[i]$ 中整除
\[
n^2+1=(n+i)(n-i).
\]
若 $p$ 在 $\ZZ[i]$ 中不可约，则 $p$ 将在 $\ZZ[i]$ 中整除 $n+i$ 或 $n-i$。在这种情形下，由于 $p$ 是实数，它也将整除其复共轭，因而整除二者的差 $2i$。这显然不可能。于是我们证明了如下结果。

\medskip
\noindent\textbf{命题 18.}
\begin{enumerate}[label=(\arabic*)]
	\item （Fermat 平方和定理）素数 $p$ 可写成两个整数平方和
	\[
	p=a^2+b^2,
	\qquad a,b\in\ZZ,
	\]
	当且仅当 $p=2$ 或 $p\equiv1\pmod4$。除交换 $a,b$ 以及改变 $a,b$ 的符号外，$p$ 表示为两个平方和的方式唯一。
	\item Gaussian 整数 $\ZZ[i]$ 中的不可约元如下：
	\begin{enumerate}[label=(\alph*)]
		\item $1+i$（它的范数为 $2$）；
		\item $\ZZ$ 中满足 $p\equiv3\pmod4$ 的素数 $p$（它们的范数为 $p^2$）；
		\item 对于 $\ZZ$ 中满足 $p\equiv1\pmod4$ 的素数 $p$，若 $p=a^2+b^2=(a+bi)(a-bi)$，则 $a+bi$ 与 $a-bi$ 是 $p$ 的两个不同不可约因子（二者范数均为 $p$）。
	\end{enumerate}
\end{enumerate}

命题 18 的第一部分是初等数论中 Fermat 的一个著名定理，它还有许多其他证明。

更一般地，整数 $n\in\ZZ$ 能否写成两个整数平方和 $n=A^2+B^2$，等价于 $n$ 是否是 Gaussian 整数中某个元素 $A+Bi$ 的范数，即
\[
n=A^2+B^2=N(A+Bi).
\]
把
\[
A+Bi=\pi_1\pi_2\cdots\pi_k
\]
唯一地（至多差单位）写成不可约元乘积，由命题 18 对 $\ZZ[i]$ 中不可约元的显式描述可知，$n$ 是一个范数当且仅当 $n$ 的那些同余于 $3$ 模 $4$ 的素因子均以偶数次幂出现。进一步地，若 $n$ 满足此条件，则 $A+Bi$ 在 $\ZZ[i]$ 中的唯一分解还允许我们数出 $n$ 表示为两个平方和的方式数，如下推论所示。

\medskip
\noindent\textbf{推论 19.}
设 $n$ 是正整数，并写
\[
n=2^k p_1^{a_1}\cdots p_r^{a_r}q_1^{b_1}\cdots q_s^{b_s},
\]
其中 $p_1,\ldots,p_r$ 是互异的同余于 $1$ 模 $4$ 的素数，$q_1,\ldots,q_s$ 是互异的同余于 $3$ 模 $4$ 的素数。则 $n$ 可写成 $\ZZ$ 中两个平方和，即
\[
n=A^2+B^2,
\qquad A,B\in\ZZ,
\]
当且仅当每个 $b_i$ 都是偶数。进一步地，若 $n$ 满足此条件，则 $n$ 表示为两个平方和的表示数为
\[
4(a_1+1)(a_2+1)\cdots(a_r+1).
\]

\begin{proof}
	推论中的第一句话已经在上面证明。现假设 $b_1,\ldots,b_s$ 都为偶数。对每个同余于 $1$ 模 $4$ 的素数 $p_i$，写
	\[
	p_i=\pi_i\overline{\pi_i},
	\]
	其中 $\pi_i$ 与 $\overline{\pi_i}$ 是命题 18(2)(c) 中的不可约元。若 $N(A+Bi)=n$，则考察范数可知，在相差单位的意义下，$A+Bi$ 在 $\ZZ[i]$ 中分解为不可约元的形式是
	\[
	A+Bi=(1+i)^k
	(\pi_1^{\alpha_{1,1}}\overline{\pi_1}^{\alpha_{1,2}})\cdots
	(\pi_r^{\alpha_{r,1}}\overline{\pi_r}^{\alpha_{r,2}})
	q_1^{b_1/2}\cdots q_s^{b_s/2},
	\]
	其中非负整数 $\alpha_{i,1},\alpha_{i,2}$ 满足
	\[
	\alpha_{i,1}+\alpha_{i,2}=a_i,
	\qquad i=1,2,\ldots,r.
	\]
	由于 $\alpha_{i,1}$ 可取 $0,1,\ldots,a_i$（而后 $\alpha_{i,2}$ 即被确定），所以在相差单位的意义下，范数为 $n$ 的 $\ZZ[i]$ 中元素 $A+Bi$ 共有
	\[
	(a_1+1)(a_2+1)\cdots(a_r+1)
	\]
	个。最后，由于 $\ZZ[i]$ 中有四个单位，推论中的第二句话随之成立。
\end{proof}

\paragraph{例}
由于
\[
493=17\cdot29
\]
且两个素数都同余于 $1$ 模 $4$，$493=A^2+B^2$ 是两个整数平方和。又因为
\[
17=(4+i)(4-i),
\qquad
29=(5+2i)(5-2i),
\]
所以 $A+Bi$ 在 $\ZZ[i]$ 中在相差单位意义下的可能分解是
\begin{align*}
(4+i)(5+2i)&=18+13i,\\
(4+i)(5-2i)&=22-3i,\\
(4-i)(5+2i)&=22+3i,\\
(4-i)(5-2i)&=18-13i.\\
\end{align*}
乘以 $-1$ 会同时改变两个符号；乘以 $i$ 会交换 $A$ 与 $B$ 并引入一个符号变化。因此
\[
493=(\pm18)^2+(\pm13)^2=(\pm22)^2+(\pm3)^2,
\]
其中所有可能的符号选择给出了 $493$ 作为两个平方和的 $16$ 种表示中的 $8$ 种；剩下的 $8$ 种通过交换两个加数得到。

类似地，整数
\[
58000957=73^2\cdot17\cdot29
\]
也恰好能以 $16$ 种方式写成两个平方和，这些表示由上面 $493=A^2+B^2$ 中的每个整数 $A,B$ 都乘以 $73$ 得到。

\subsection*{小结}

综上，我们在含单位交换环的各类之间有如下包含关系：
\[
\text{域}\subset \text{Euclidean整环}\subset \text{P.I.D.s}\subset \text{U.F.D.s}\subset \text{整环},
\]
且所有包含都是真包含。回忆，$\ZZ$ 是不是域的 Euclidean整环；二次整数环
\[
\ZZ\left[{1+\sqrt{-19}\over2}\right]
\]
是主理想整环但不是 Euclidean整环；$\ZZ[x]$ 是唯一分解整环（第 9 章定理 7）但不是主理想整环；而 $\ZZ[\sqrt{-5}]$ 是整环但不是唯一分解整环。

\paragraph{练习}
\begin{enumerate}[label=\arabic*.]
	\item 令 $G=\QQ^\times$ 为非零有理数的乘法群。若 $\alpha=p/q\in G$，其中 $p,q$ 是互素整数，定义映射
	\[
	\varphi:G\longrightarrow G
	\]
	为在 $p$ 与 $q$ 的素数幂分解中互换素数 $2$ 与 $3$（例如
	\[
	\varphi(2^4 3^1 5^1 13^2)=3^4 2^1 5^1 13^2,
	\qquad
	\varphi(3/16)=\varphi(3/2^4)=2/3^4=2/81,
	\]
	并且 $\varphi$ 在所有分子、分母都与 $2$ 和 $3$ 互素的有理数上为恒等）。
	\begin{enumerate}[label=(\alph*)]
		\item 证明 $\varphi$ 是群同构。
		\item 证明群 $G$ 到自身有无穷多个同构。
		\item 证明上述同构都不能扩张为环 $\QQ$ 到自身的同构。事实上，证明 $\QQ$ 的唯一环同构是恒等映射。
	\end{enumerate}
	
	\item 设 $a,b$ 是唯一分解整环 $R$ 中的非零元素。证明 $a,b$ 有最小公倍数（参见第 1 节练习 11），并像命题 13 描述最大公因子那样，用 $a,b$ 的素因子分解描述它。
	
	\item 求整数
	\[
	2130797=17^2\cdot73\cdot101
	\]
	表示为两个平方和的全部表示。
	
	\item 证明：若一个整数是两个有理数平方的和，则它也是两个整数平方的和（例如
	\[
	13=(1/5)^2+(18/5)^2=2^2+3^2
	\]
	）。
	
	\item 令 $R=\ZZ[\sqrt n]$，其中 $n$ 是大于 $3$ 的无平方因子整数。
	\begin{enumerate}[label=(\alph*)]
		\item 证明 $2,\sqrt n$ 与 $1+\sqrt n$ 在 $R$ 中不可约。
		\item 证明 $R$ 不是 U.F.D.。推出当 $D\equiv2,3\pmod4$ 且 $D<-3$ 时，二次整数环 $\mcO$ 不是 U.F.D.（因此也不是 Euclidean整环，也不是 P.I.D.）。\emph{提示：}说明 $\sqrt n$ 或 $1+\sqrt n$ 不是素元。
		\item 给出 $R$ 中一个非主理想的显式例子。\emph{提示：}利用 (b)，考虑包含非素理想 $(\sqrt n)$ 或 $(1+\sqrt n)$ 的极大理想。
	\end{enumerate}
	
	\item
	\begin{enumerate}[label=(\alph*)]
		\item 证明商环 $\ZZ[i]/(1+i)$ 是一个 $2$ 元域。
		\item 设 $q\in\ZZ$ 是满足 $q\equiv3\pmod4$ 的素数。证明商环 $\ZZ[i]/(q)$ 是一个含 $q^2$ 个元素的域。
		\item 设 $p\in\ZZ$ 是满足 $p\equiv1\pmod4$ 的素数，并按命题 18 写成 $p=\pi\overline\pi$。证明中国剩余定理（第 7.6 节定理 17）的假设满足，并且有环同构
		\[
		\ZZ[i]/(p)\cong \ZZ[i]/(\pi)\times \ZZ[i]/(\overline\pi)
		\]
		作为环。证明商环 $\ZZ[i]/(p)$ 的阶为 $p^2$，并推出 $\ZZ[i]/(\pi)$ 与 $\ZZ[i]/(\overline\pi)$ 都是 $p$ 元域。
	\end{enumerate}
	
	\item 设 $\pi$ 是 $\ZZ[i]$ 中的不可约元。
	\begin{enumerate}[label=(\alph*)]
		\item 对任意整数 $n\geq0$，证明 $(\pi^{n+1})=\pi^{n+1}\ZZ[i]$ 是 $(\pi^n)=\pi^n\ZZ[i]$ 中的一个理想，并且乘以 $\pi^n$ 诱导出加法 Abel 群同构
		\[
		\ZZ[i]/(\pi)\cong (\pi^n)/(\pi^{n+1}).
		\]
		\item 证明
		\[
		|\ZZ[i]/(\pi^n)|=|\ZZ[i]/(\pi)|^n.
		\]
		\item 证明：对 $\ZZ[i]$ 中任意非零元素 $\alpha$，商环 $\ZZ[i]/(\alpha)$ 的阶等于 $N(\alpha)$。\emph{提示：}结合 (b)、中国剩余定理以及上一题的结果。
	\end{enumerate}
	
	\item 令 $R$ 为二次整数环 $\ZZ[\sqrt{-5}]$，并定义理想
	\[
	I_2=(2,1+\sqrt{-5}),\qquad
	I_3=(3,2+\sqrt{-5}),\qquad
	I_3'=(3,2-\sqrt{-5}).
	\]
	\begin{enumerate}[label=(\alph*)]
		\item 证明 $2,3,1+\sqrt{-5},1-\sqrt{-5}$ 在 $R$ 中都是不可约元，且任意两个都不相伴，并且
		\[
		6=2\cdot3=(1+\sqrt{-5})(1-\sqrt{-5})
		\]
		是 $6$ 在 $R$ 中分解为不可约元的两个不同分解。
		\item 证明 $I_2,I_3,I_3'$ 是 $R$ 中的素理想。\emph{提示：}对 $I_3$，注意由环的第三同构定理有
		\[
		R/I_3\cong (R/(3))/(I_3/(3)).
		\]
		证明 $R/(3)$ 有 $9$ 个元素，$I_3/(3)$ 有 $3$ 个元素，且 $R/I_3$ 作为加法 Abel 群同构于 $\ZZ/3\ZZ$。推出 $I_3$ 是极大理想（因此是素理想），并且 $R/I_3\cong \ZZ/3\ZZ$ 作为环。
		\item 说明 (a) 中的分解给出理想等式
		\[
		(6)=(2)(3),\qquad
		(6)=(1+\sqrt{-5})(1-\sqrt{-5}).
		\]
		证明这两个理想分解给出了理想 $(6)$ 分解为素理想乘积的同一个分解（参见第 2 节练习 5）。
	\end{enumerate}
	
	\item 假设二次整数环 $\mcO$ 是 P.I.D.。证明 $\mcO$ 上的域范数绝对值 $|N|$（参见第 7.1 节）是 $\mcO$ 上的 Dedekind--Hasse范数。推出：若二次整数环 $\mcO$ 存在任何 Dedekind--Hasse范数，则事实上 $\mcO$ 上的域范数绝对值已经给出了一个 Dedekind--Hasse范数。\emph{提示：}若 $\alpha,\beta\in\mcO$，则 $(\alpha,\beta)=(\gamma)$ 对某个 $\gamma\in\mcO$ 成立。证明若 $\beta$ 不整除 $\alpha$，则
	\[
	0<|N(\gamma)|<|N(\beta)|
	\]
	——使用 $\mcO$ 中的单位恰为范数为 $\pm1$ 的元素这一事实。
	
	\noindent\emph{说明：}如果 $\mcO$ 关于某个范数是 Euclidean整环，并不一定意味着它关于域范数绝对值是 Euclidean整环（虽然对 $D<0$ 这是真的，参见第 1 节练习 8）。一个例子是 $D=69$（参见 D. Clark, \emph{A quadratic field which is Euclidean but not norm-Euclidean}, Manuscripta Math. 83 (1994), pp. 327--330）。
	
	\item （$k$ 阶 Euclidean整环）设 $R$ 是整环，$N:R\to\ZZ^+\cup\{0\}$ 是 $R$ 上的范数。若对任意 $a,b\in R$，$b\neq0$，存在 $q,r\in R$ 满足
	\[
	a=qb+r,
	\qquad r=0\text{ 或 }N(r)<N(b),
	\]
	则 $R$ 关于 $N$ 是 Euclidean的。现在弱化这个条件：假设对任意 $a,b\in R$，$b\neq0$，存在 $q,q',r,r'\in R$ 使得
	\[
	a=qb+r,
	\qquad
	b=q'r+r',
	\qquad r'=0\text{ 或 }N(r')<N(b),
	\]
	即两次除法之后的余数变小。称这样的整环为 \emph{2 阶 Euclidean整环}。
	\begin{enumerate}[label=(\alph*)]
		\item 证明：在 2 阶 Euclidean整环中迭代这些除法会产生 $a$ 与 $b$ 的一个最大公因子，并且该最大公因子是 $a$ 与 $b$ 的线性组合。推出：2 阶 Euclidean整环中的每个有限生成理想都是主理想。（不过存在不是 P.I.D. 的 2 阶 Euclidean整环。）\emph{提示：}模仿定理 4 的证明。
		\item 证明：若一个 2 阶 Euclidean整环中每个非零非单位都可以分解为有限个不可约元的乘积，则它是唯一分解整环。\emph{提示：}先证明不可约元都是素元，如下：若 $p$ 不可约，$p\mid ab$，且 $p$ 不整除 $a$，利用 (a) 写出
		\[
		px+ay=1
		\]
		对某些 $x,y$ 成立。两边乘以 $b$，推出 $p\mid b$，所以 $p$ 是素元。然后沿用定理 14 中唯一性的证明。
		\item 对任意整数 $k\geq1$，作出显然的推广以定义 $k$ 阶 Euclidean整环。证明若把“2 阶 Euclidean”替换为“$k$ 阶 Euclidean”，(a) 与 (b) 仍成立。
	\end{enumerate}
	
	\noindent\emph{说明：}存在 2 阶 Euclidean 但不是 Euclidean的环。也存在关于给定范数不是 Euclidean、但关于该范数是 $k$ 阶 Euclidean的环（例如环 $\ZZ[\sqrt{14}]$ 关于通常范数 $N(a+b\sqrt{14})=|a^2-14b^2|$ 不是 Euclidean的，但关于此范数是 2 阶 Euclidean的）。$k$ 阶 Euclidean 条件还与如下问题有关：由 $R$ 中元素构成的可逆 $n\times n$ 矩阵群 $\operatorname{GL}_n(R)$ 是否由初等矩阵生成（初等矩阵指主对角线上全为 $1$，在主对角线外某一位置有一个 $1$，其余位置全为 $0$ 的矩阵）。
	
	\item （P.I.D. 的刻画）证明：$R$ 是 P.I.D. 当且仅当 $R$ 是同时为 Bezout整环的 U.F.D.（参见第 2 节练习 7）。\emph{提示：}一个方向由定理 14 给出。反过来，令 $a$ 是理想 $I$ 中一个非零元素，使其不可约因子个数最少。通过证明若存在 $b\in I$ 且 $b\notin(a)$，则 $(a,b)=(d)$ 会导致矛盾，从而证明 $I=(a)$。
\end{enumerate}
